\chapter{Short-Time Existence}

\section*{1. Variation formulas}

\begin{lemma}
\label{ch03-lem-metric-inverse-evolves}
\dcref{Lemma 3.1}
\group{variation-formulas}
\source{chow-knopf-ricci-flow:page-0080}
If \(\partial_t g_{ij}=h_{ij}\), then the inverse metric satisfies
\[
\partial_t g^{ij}=-g^{ik}g^{j\ell}h_{k\ell}.
\]
\end{lemma}

\begin{proof}
\source{chow-knopf-ricci-flow:page-0080, chow-knopf-ricci-flow:page-0081}
Differentiating \(\delta^i_{\ell}=g^{ik}g_{k\ell}\) gives
\[
0=(\partial_t g^{ik})g_{k\ell}+g^{ik}h_{k\ell},
\]
and multiplying by \(g^{j\ell}\) yields the stated formula.
\end{proof}

\begin{lemma}
\label{ch03-lem-connection-evolves}
\dcref{Lemma 3.2}
\group{variation-formulas}
\source{chow-knopf-ricci-flow:page-0081}
If \(\partial_t g_{ij}=h_{ij}\), then the Levi-Civita connection varies by
\[
\partial_t\Gamma^k_{ij}
=\frac12 g^{k\ell}\bigl(\nabla_i h_{j\ell}+\nabla_j h_{i\ell}-\nabla_\ell h_{ij}\bigr).
\]
\end{lemma}

\begin{proof}
\source{chow-knopf-ricci-flow:page-0081}
In geodesic coordinates at a point, \(\Gamma^k_{ij}=0\) and \(\partial_i=\nabla_i\) at that point. Differentiating the coordinate formula for \(\Gamma^k_{ij}\) and then invoking tensoriality gives the result globally.
\end{proof}

\begin{lemma}
\label{ch03-lem-riemann-evolves}
\dcref{Lemma 3.3}
\group{variation-formulas}
\source{chow-knopf-ricci-flow:page-0081,chow-knopf-ricci-flow:page-0082}
If \(\partial_t g_{ij}=h_{ij}\), then the Riemann tensor varies by
\[
\partial_t R^\ell{}_{ijk}
=\frac12 g^{\ell p}\!\left(
\nabla_i\nabla_j h_{kp}+\nabla_i\nabla_k h_{jp}-\nabla_i\nabla_p h_{jk}
-\nabla_j\nabla_i h_{kp}-\nabla_j\nabla_k h_{ip}+\nabla_j\nabla_p h_{ik}
\right).
\]
\end{lemma}

\begin{proof}
\source{chow-knopf-ricci-flow:page-0081,chow-knopf-ricci-flow:page-0082}
Differentiate the coordinate expression for \(R^\ell{}_{ijk}\), use the variation formula for \(\Gamma^k_{ij}\), and reduce at a point using geodesic coordinates. The result is tensorial, so the pointwise computation is global.
\end{proof}

\begin{remark}
\label{ch03-rem-riemann-alt-form}
\dcref{Remark 3.4}
\group{variation-formulas}
\source{chow-knopf-ricci-flow:page-0082}
Commuting derivatives gives an alternative form for \(\partial_t R^\ell{}_{ijk}\) in which curvature terms appear explicitly.
\end{remark}

\begin{lemma}
\label{ch03-lem-ricci-evolves}
\dcref{Lemma 3.5}
\group{variation-formulas}
\source{chow-knopf-ricci-flow:page-0082}
If \(\partial_t g_{ij}=h_{ij}\), then the Ricci tensor varies by
\[
\partial_t R_{jk}
=\frac12 g^{pq}\bigl(\nabla_q\nabla_j h_{kp}
+\nabla_q\nabla_k h_{jp}
-\nabla_q\nabla_p h_{jk}
-\nabla_j\nabla_k h_{qp}\bigr).
\]
\end{lemma}

\begin{proof}
\source{chow-knopf-ricci-flow:page-0082}
This is the contraction \(i=\ell\) of Lemma~\ref{ch03-lem-riemann-evolves}.
\end{proof}

\begin{remark}
\label{ch03-rem-lichnerowicz-laplacian}
\dcref{Remark 3.6}
\group{variation-formulas}
\source{chow-knopf-ricci-flow:page-0082}
Writing \((\delta h)_k=-g^{ij}\nabla_i h_{jk}\), and denoting by \(\Delta_L\) the Lichnerowicz Laplacian, the variation of \(\Rc\) can be rewritten in terms of \(\Delta_L h\), \(\nabla\nabla H\), and \(\delta h\), where \(H=\tr_g h\).
\end{remark}

\begin{lemma}
\label{ch03-lem-scalar-curvature-evolves}
\dcref{Lemma 3.7}
\group{variation-formulas}
\source{chow-knopf-ricci-flow:page-0082}
If \(\partial_t g_{ij}=h_{ij}\), then the scalar curvature varies by
\[
\partial_t R
=-\Delta H+\nabla^p\nabla^q h_{pq}-\langle h,\Rc\rangle.
\]
\end{lemma}

\begin{proof}
\source{chow-knopf-ricci-flow:page-0082}
Combine Lemmas~\ref{ch03-lem-metric-inverse-evolves} and \ref{ch03-lem-ricci-evolves}, then contract the Ricci variation and simplify.
\end{proof}

\begin{remark}
\label{ch03-rem-scalar-curvature-invariant-form}
\dcref{Remark 3.8}
\group{variation-formulas}
\source{chow-knopf-ricci-flow:page-0082}
The scalar curvature variation can also be written invariantly as
\[
\partial_t R=-\Delta H+\operatorname{div}(\operatorname{div} h)-\langle h,\Rc\rangle.
\]
\end{remark}

\begin{lemma}
\label{ch03-lem-volume-element-evolves}
\dcref{Lemma 3.9}
\group{variation-formulas}
\source{chow-knopf-ricci-flow:page-0083}
The volume form satisfies
\[
\partial_t d\mu=\frac12(g^{ij}\partial_t g_{ij})\,d\mu=\frac{H}{2}\,d\mu.
\]
\end{lemma}

\begin{proof}
\source{chow-knopf-ricci-flow:page-0083}
Differentiate \(d\mu=\sqrt{\det g}\,dx^1\wedge\cdots\wedge dx^n\) and use \(\partial_t\log\det g=g^{ij}\partial_t g_{ij}\).
\end{proof}

\begin{corollary}
\label{ch03-cor-total-scalar-curvature-evolves}
\dcref{Corollary 3.10}
\group{variation-formulas}
\source{chow-knopf-ricci-flow:page-0083}
The total scalar curvature evolves by
\[
\frac{d}{dt}\int_{\M^n}R\,d\mu
=\int_{\M^n}\Bigl(\frac12 RH-\langle \Rc,h\rangle\Bigr)\,d\mu.
\]
\end{corollary}

\begin{lemma}
\label{ch03-lem-length-evolves}
\dcref{Lemma 3.11}
\group{variation-formulas}
\source{chow-knopf-ricci-flow:page-0083,chow-knopf-ricci-flow:page-0084}
If \(\gamma_t\) is a time-dependent family of curves with fixed endpoints and \(\partial_t g=h\), then
\[
\frac{d}{dt}L_t(\gamma_t)
=\frac12\int_{\gamma_t}h(S,S)\,ds-\int_{\gamma_t}\langle\nabla_S S,V\rangle\,ds,
\]
where \(S\) is the unit tangent and \(V\) is the variation vector field.
\end{lemma}

\begin{proof}
\source{chow-knopf-ricci-flow:page-0083,chow-knopf-ricci-flow:page-0084}
Differentiate the length integral, use \([U,V]=0\), and integrate by parts the term \(\langle S,\nabla_S V\rangle\); the endpoint contribution vanishes because the endpoints are fixed.
\end{proof}

\section*{2. The linearization of Ricci}

\begin{remark}
\label{ch03-rem-symbol-of-composition}
\dcref{Remark 3.12}
\group{linearization-of-ricci}
\source{chow-knopf-ricci-flow:page-0085, chow-knopf-ricci-flow:page-0086, chow-knopf-ricci-flow:page-0087, chow-knopf-ricci-flow:page-0088, chow-knopf-ricci-flow:page-0089}
The diffeomorphism invariance of \(\Rc\) and \(\Rm\) implies the contracted second Bianchi identity, the first and second Bianchi identities, and the algebraic identities governing the symbols of \(\delta_g^*\), \(D(\Rc_g)\), and \(B_g\).
\end{remark}

\section*{3. Short-Time Existence}

\begin{theorem}
\label{ch03-thm-short-time-existence}
\dcref{Theorem 3.13}
\group{short-time-existence}
\source{chow-knopf-ricci-flow:page-0091}
If \((\M^n,g_0)\) is a closed Riemannian manifold, there exists a unique solution \(g(t)\) to the Ricci flow on some interval \([0,\varepsilon)\) with \(g(0)=g_0\).
\end{theorem}

\begin{remark}
\label{ch03-rem-lifetime-lower-bound}
\dcref{Remark 3.14}
\group{short-time-existence}
\source{chow-knopf-ricci-flow:page-0092}
The maximal lifespan of the Ricci flow solution is bounded below by a universal constant times the reciprocal of the initial curvature maximum.
\end{remark}

\begin{lemma}
\label{ch03-lem-deTurck-flow-diffeomorphisms}
\dcref{Lemma 3.15}
\group{short-time-existence}
\source{chow-knopf-ricci-flow:page-0095}
If \(\{X_t\}\) is a continuous time-dependent family of vector fields on a compact manifold, then the ODE \(\partial_t\varphi_t=X_t\circ\varphi_t\), \(\varphi_0=\mathrm{id}\), has a solution by diffeomorphisms on the same time interval.
\end{lemma}

\begin{proof}
\source{chow-knopf-ricci-flow:page-0095}
The flow exists locally in coordinates as an ODE in \(\mathbb R^n\); compactness gives a uniform existence time, and iteration extends the flow on the whole interval.
\end{proof}

\begin{remark}
\label{ch03-rem-noncompact-flow-failure}
\dcref{Remark 3.16}
\group{short-time-existence}
\source{chow-knopf-ricci-flow:page-0095}
The lemma may fail on noncompact manifolds; on \(\mathbb R\), the ODE \(u'=u^2\) provides a counterexample.
\end{remark}

\begin{lemma}
\label{ch03-lem-riemann-tensor-decomposition}
\dcref{Lemma 3.17}
\group{deTurck-formulation}
\source{chow-knopf-ricci-flow:page-0096}
The variation of the Ricci tensor may be written as
\[
D(\Rc_g)(h)=-\frac12\Delta_L h-\delta_g^*(\delta_g G(h)),
\]
where \(G(h)=h-\frac12(\tr_g h)\,g\).
\end{lemma}

\begin{proof}
\source{chow-knopf-ricci-flow:page-0096}
Rearrange the standard linearization of \(\Rc\) into its Lichnerowicz part and the divergence term involving the Einstein operator \(G\).
\end{proof}

\begin{lemma}
\label{ch03-lem-harmonic-map-laplacian}
\dcref{Lemma 3.18}
\group{harmonic-map-flow}
\source{chow-knopf-ricci-flow:page-0097, chow-knopf-ricci-flow:page-0098, chow-knopf-ricci-flow:page-0099}
If \(f:(\M^n,g)\to(\N^n,h)\) is a diffeomorphism, then its harmonic map Laplacian satisfies
\[
(\Delta_{g,h}f)^\gamma(x)
=\bigl[(f^{-1})^*g\bigr]^{\alpha\beta}
\left(-\Gamma\bigl((f^{-1})^*g\bigr)^\gamma_{\alpha\beta}+\Gamma(h)^\gamma_{\alpha\beta}\right)\!(f(x)).
\]
\end{lemma}

\begin{proof}
\source{chow-knopf-ricci-flow:page-0099}
Compute the Christoffel symbols of a pullback metric, then rewrite the harmonic map Laplacian in terms of those symbols by substituting \(\varphi=f\) and \(\kappa=(f^{-1})^*g\).
\end{proof}

\begin{corollary}
\label{ch03-cor-identity-harmonic-map}
\dcref{Corollary 3.19}
\group{harmonic-map-flow}
\source{chow-knopf-ricci-flow:page-0099}
For a diffeomorphism \(f\) and \(h=(f^{-1})^*g\), the identity map from \((\M^n,g)\) to \((\M^n,h)\) has harmonic map Laplacian given by the connection difference \(g^{\alpha\beta}\bigl(-\Gamma(g)^\gamma_{\alpha\beta}+\Gamma(h)^\gamma_{\alpha\beta}\bigr)\).
\end{corollary}

\begin{corollary}
\label{ch03-cor-positive-ricci-identity-harmonic}
\dcref{Corollary 3.20}
\group{harmonic-map-flow}
\source{chow-knopf-ricci-flow:page-0099, chow-knopf-ricci-flow:page-0100}
If \((\M^n,g)\) has strictly positive Ricci curvature, then \(\mathrm{id}:(\M^n,g)\to(\M^n,\Rc(g))\) is a harmonic map.
\end{corollary}

\begin{remark}
\label{ch03-rem-negative-ricci-target}
\dcref{Remark 3.21}
\group{harmonic-map-flow}
\source{chow-knopf-ricci-flow:page-0100}
The same harmonic-map conclusion holds with target metric \(-\Rc(g)\) when \(\Rc<0\).
\end{remark}

\begin{remark}
\label{ch03-prob-cross-curvature-harmonic-map}
\dcref{Problem 3.22}
\group{cross-curvature-flow}
\source{chow-knopf-ricci-flow:page-0100}
Given a Riemannian manifold \((\M^n,g)\), determine whether there exists a metric \(\gamma\neq g\) such that \(\mathrm{id}:(\M^n,\gamma)\to(\M^n,g)\) is harmonic.
\end{remark}

\begin{remark}
\label{ch03-ex-cross-curvature-harmonic-map}
\dcref{Exercise 3.23}
\group{cross-curvature-flow}
\source{chow-knopf-ricci-flow:page-0101}
Verify the divergence identity for the cross curvature tensor and use it to show that if the sectional curvatures of \((\M^3,g)\) are everywhere positive or everywhere negative, then \(c_g\) is a metric and \(\mathrm{id}:(\M^3,c_g)\to(\M^3,g)\) is harmonic.
\end{remark}

\begin{proposition}
\label{ch03-prop-cross-curvature-monotonicity}
\dcref{Proposition 3.24}
\group{cross-curvature-flow}
\source{chow-knopf-ricci-flow:page-0101}
Along a solution of the cross curvature flow, \(\Vol(E)\) is nondecreasing and the integral of \(\frac13(g^{ij}E_{ij})-(\det E/\det g)^{1/3}\) is nonincreasing.
\end{proposition}

\begin{conjecture}
\label{ch03-conj-cross-curvature-global-existence}
\dcref{Conjecture 3.25}
\group{cross-curvature-flow}
\source{chow-knopf-ricci-flow:page-0101}
If \((\M^3,g_0)\) has negative sectional curvature, the cross curvature flow should exist for all positive time and converge.
\end{conjecture}

\begin{remark}
\label{ch03-rem-andrews-observation}
\dcref{Remark 3.26}
\group{cross-curvature-flow}
\source{chow-knopf-ricci-flow:page-0102}
Andrews observed that certain hypersurface Gauss curvature flows induce cross curvature flow on the underlying metric.
\end{remark}

\begin{lemma}
\label{ch03-lem-deTurck-diffeomorphisms-harmonic-map}
\dcref{Lemma 3.27}
\group{harmonic-map-flow}
\source{chow-knopf-ricci-flow:page-0102, chow-knopf-ricci-flow:page-0103}
The DeTurck diffeomorphisms \(\varphi_t\) solving \(\partial_t\varphi_t=-W\circ\varphi_t\) satisfy the harmonic map heat flow with domain metric \(\bar g(t)\) and codomain metric \(\tilde g\).
\end{lemma}

\begin{proof}
\source{chow-knopf-ricci-flow:page-0102, chow-knopf-ricci-flow:page-0103}
Substitute the DeTurck vector field into the harmonic map Laplacian formula and identify the resulting evolution equation with \(\partial_t\varphi_t\).
\end{proof}

\begin{definition}
\label{ch03-def-harmonic-coordinates}
\dcref{Definition 3.28}
\group{harmonic-coordinates}
\source{chow-knopf-ricci-flow:page-0103}
Local coordinates \(\{x^i\}\) are harmonic if each coordinate function is harmonic, equivalently if \(g^{jk}\Gamma^i_{jk}=0\).
\end{definition}

\begin{theorem}
\label{ch03-thm-local-elliptic-solvability}
\dcref{Theorem 3.29}
\group{harmonic-coordinates}
\source{chow-knopf-ricci-flow:page-0104}
Let \(F(x,u,\partial u,\ldots,\partial^ku)\) be smooth. If \(F\) is elliptic at a jet \(v\) near \(x_0\) and \(F(x_0,v,\partial v,\ldots,\partial^kv)=0\), then there is a local solution \(u\) near \(x_0\) matching the derivatives of \(v\) up to order \(k-1\) at \(x_0\).
\end{theorem}

\begin{corollary}
\label{ch03-cor-harmonic-coordinates-exist}
\dcref{Corollary 3.30}
\group{harmonic-coordinates}
\source{chow-knopf-ricci-flow:page-0104}
Every point of a smooth manifold admits harmonic coordinates in a neighborhood.
\end{corollary}

\begin{lemma}
\label{ch03-lem-harmonic-coordinate-regularity}
\dcref{Lemma 3.31}
\group{harmonic-coordinates}
\source{chow-knopf-ricci-flow:page-0104}
If a metric or tensor is \(C^{k,\alpha}\) or \(C^\infty\) in one chart, then it has the same regularity in harmonic coordinates.
\end{lemma}

\begin{lemma}
\label{ch03-lem-ricci-in-harmonic-coordinates}
\dcref{Lemma 3.32}
\group{harmonic-coordinates}
\source{chow-knopf-ricci-flow:page-0105}
In harmonic coordinates, the Ricci tensor satisfies
\[
-2R_{ij}=\Delta(g_{ij})+Q_{ij}(g^{-1},\partial g),
\]
where \(Q\) is quadratic in \(g^{-1}\) and \(\partial g\).
\end{lemma}

\begin{corollary}
\label{ch03-cor-ricci-flow-heat-equation}
\dcref{Corollary 3.33}
\group{harmonic-coordinates}
\source{chow-knopf-ricci-flow:page-0105}
In harmonic coordinates, the Ricci flow becomes a nonlinear heat equation for the metric components:
\[
\partial_t g_{ij}=\Delta(g_{ij})+Q_{ij}(g^{-1},\partial g).
\]
\end{corollary}

\begin{remark}
\label{ch03-rem-harmonic-coordinate-caution}
\dcref{Remark 3.34}
\group{harmonic-coordinates}
\source{chow-knopf-ricci-flow:page-0105}
The heat-equation form of the Ricci flow is not tensorial, and coordinates that are harmonic at one time need not remain harmonic later.
\end{remark}
